\subsection{Simultaneous activation and completability}
\begin{proposicion}[Norm of the error on the activation cube]\label{prop:norma-llegada}
Let \(Z(Y)\) be the probability of avoiding the active pairs of \(Y\) under a uniform rook placement on
\(A\).
Let \(\epsilon_A\) be the maximum diagonal load of \(p^{\rm T}\). There is a class threshold \(\eta>0\)
such that if \(\rho\ge2\) and \(\rho^2(\epsilon_A+1/j)\le\eta\), then, for large \(j\), \(Z(Y)>0\) on
every face and \(e=\log Z-\Lambda\) satisfies
\[
 \sum_{U\subseteq A}\rho^{|U|}|\Delta_Ue|\le C_\rho j^2.
\]
\end{proposicion}
\begin{proof}
We use the finite algebra of functions on the cube with \(z^Uz^V=z^{U\cup V}\), where
\(z^U(Y)=\mathbf1_{\{U\subseteq Y\}}\). The norm \(\|\sum a_Uz^U\|_\rho
=\sum|a_U|\rho^{|U|}\) is submultiplicative and dominates the absolute value of every evaluation. Let
\(a_U=\Delta_U\mathbf1_{\{\text{no diagonal pairs}\}}\). If \(p_U\) is the probability that the rooks
occupy \(U\), inversion on their intersection with \(Y\) gives \(Z=\sum_Ua_Up_U z^U\). Axial
incompatibilities have \(p_U=0\).

For a compatible prescription of \(h\) sites let \(T_U=\sum_{(i,l)\in U}(1,\delta_i,\delta_l)\), of total
variation norm \(3h\). The balancing and gamma truncation of Proposition~\ref{prop:menores-llegada} allow
\(h\le a_*j\) with a small class constant \(a_*\), keeping the same \(Q_*\) for parent and minor. On the
ball \(\|T\|\le c j\) we balance \(jD\) in the surviving measures normalized by \(1-h/j\). The baseline
Jacobian remains invertible by perturbation. If \(u_i,v_l\) are the potentials of this balancing, the
change of scale of the permanent adds to \(\log(m!/j!)+h\log j\), with \(m=j-h\), the negative sum of the
potentials on the surviving margins. Its linear term defines weights \(b_i\asymp D_i\). The remaining
term, together with the difference between \(\log H_{Q_*}\) for the two scalings, defines a field
\(L_*(T)\) that vanishes at the origin. We extend the factorial term via \(\log\Gamma\). On the ball, the
sum of potentials and the factorial cost \(O(j)\), and the logarithm of the truncation costs \(O(1)\). By
Cauchy's estimate in two directions, \(\|D^2L_*\|\le C/j\) on the ball of half the radius. Its derivative
at zero is \(O(1/j)\) by centering and \(\log j-(\log\Gamma)'(j+1)=O(1/j)\). Thus
\[
 p_U=\Big(\prod_{i\in U}b_i\Big)e^{L_*(T_U)}(1+O(e^{-cj})).
\]
For \(h>a_*j\), \(\operatorname{per}D_{-U}\le1\) by the row masses, and \(\operatorname{per}D\ge
j!/j^j\ge e^{-j}\) by the van der Waerden inequality. On both intervals we have \(p_U\le C_0^h\prod
b_i\).

The weight \(w_U=a_Uz^U\prod b_i\) factorizes over the components of the diagonal graph induced on \(U\),
since the avoidance indicator factorizes before differences are taken. An isolated site makes the weight
vanish. Two components of an admissible configuration share neither a site nor an axial fiber and have no
diagonal attack between them. The difference formula gives \(|a_U|\le2^{|U|}\). Every connected component
can be traversed along a spanning tree. For a tree of size \(v\) rooted at \(i\), summing from the leaves
costs \(C\epsilon_A\) per neighbor. Including the norm, the \(v^{v-2}\) choices of labeled tree, and a
fixed reserve \(e^{Bv}\), the mass is at most
\[
 b_i(C_B\rho)^v\,v^{v-2}(C\epsilon_A)^{v-1}/(v-1)! .
\]
Since \(\sum b_i=O(j)\), the total mass is \(O(S_*)\), \(S_*=j\rho^2(\epsilon_A+1/j)\). Objects
incompatible with a component of \(v\) sites have weight \(O(vS_*/j)\), by summing over its rows,
columns, and diagonal neighbors.

For a small linear functional \(\xi\) we tilt by \(e^{\xi[T_U]}\). By the exponential formula, the
logarithm of the gas of compatible components is the sum over families whose incompatibility graph is
connected. For a fixed edge order, the connected subgraphs with the same minimum spanning tree form a
Boolean interval. We may add any edge that comes after every edge on the tree path between its endpoints.
Their alternating sum has absolute value zero or one, so the trees bound the coefficient. Summing the
children of a root of size \(v\) recursively gives an exponential factor at most
\(\exp(\sum_{\gamma'\not\sim\gamma}|w_{\gamma'}|e^{|\gamma'|})\le e^v\). The criterion holds for small
\(S_*/j\) with the reserve above. This gives \(\|\log Z_\xi\|=O(S_*)\). One and two marks cost the size
factors \(3v,(3v)^2\), absorbed by the reserve. The tilted mean \(m_\xi\) of \(T_U\) and its derivative
have norm \(O(S_*)\). By the Hessian estimate above, the equation \(\xi=DL_*(m_\xi)\) is contractive with
constant \(CS_*/j<1/2\). Iteration can start from zero, so the solution is real at every Boolean
evaluation.

In \(D_\xi=Z_\xi\exp(L_*(m_\xi)-\xi[m_\xi])\) both exponentials are real and positive on each face.
Centering \(T_U\) at \(m_\xi\) cancels the linear part. The logarithm of the centered moment series has
norm \(\le CS_* t^2\) when marking with a small radius \(t\), since two or more marks cost \(C t^2v^2
e^{3tv}\) per connected family. By Cauchy's estimate its \(\ell\)-th moment has projective norm
\(\le\ell!t^{-\ell}\exp(CS_*t^2)\le[C(\sqrt{\ell S_*}+\ell)]^\ell\), taking
\(t=\min(t_0,\sqrt{\ell/(1+S_*)})\). For \(\exp(L_*(m_\xi+y)-L_*(m_\xi)-\xi[y])\) the bound is
\(\exp(C\|y\|^2/j)\). Choosing a radius proportional to \(\sqrt{j\ell}\) bounds its coefficients by
\((C/\sqrt{j\ell})^\ell\) up to \(\ell\le\kappa j\). From degree two onward, separating the first two
factors and summing with ratio \(C(\sqrt{\eta}+\sqrt\kappa)<1\) gives cost \(C(1+S_*)/j\). The projective
norm allows the contraction without a dimension factor.

For \(|U|\le a'j\) we truncate this Taylor expansion at \(\lfloor\kappa j\rfloor\) within the radius
\(c\sqrt\kappa\,j\), with \(\|T_U-m_\xi\|\le3a'j+CS_*\). Choosing \(a'\ll\sqrt\kappa\) and then \(\eta\)
leaves remainder \(O(e^{-c'\kappa j})\). For \(|U|>a'j\), the exponential reserve and the bound
\(C_0^{|U|}\) give \(e^{-a'j+CS_*}\), also for the truncated polynomial since
\(\ell!\|[y^\ell]\|\le(C\sqrt\kappa)^\ell\). The normalization factor \(e^{CS_*}\) is absorbed by
choosing \(\eta\) after the three positive rates. For large \(j\) this gives
\begin{equation}\label{eq:positividad-cubo}
 \|D_\xi^{-1}Z-1\|_\rho\le C(1+S_*)/j+Ce^{-c'j}<1/2 .
\end{equation}
The series for \(\log(1+w)\) converges in norm. Thus \(Z(Y)>0\) and \(\|\log Z\|_\rho\le C(1+S_*)\).

For the predictor, we balance \(p^{\rm T}_e e^{v_e}\) simultaneously in the same algebra. The baseline
margin inverse has bounded absolute row sums. By the gap, its solution for bounded data has Euclidean
norm \(O(\sqrt j)\), and the equation in each coordinate uses a section of norm \(O(j^{-1/2})\). For
\(h_f=1-\sum_{e\in f}z_ep^{\rm T}_e e^{v_e}\), the norm of \(1-h_f\) is \(\le \rho e^{\|v\|}\epsilon_A\).
The local odds terms contain \(\log(1-p^{\rm T}_e e^{v_e})\). The resulting field and its Lipschitz
constant are \(O(\rho^2\epsilon_A+\rho/j)\). The balanced iteration gives \(\|e^{v_e}-1\|_\rho=O(\eta)\).
The resulting real saddle is the maximum. For \(d=2,4\) the tangent to \(-t\log t+(d-1)(1-t)\log(1-t)\)
at \(q<e^{-4}\) dominates the function on \([0,1/d]\) by concavity and on \([1/d,1]\) by the two
endpoints. At one its excess is \(-\log q-d(1-q)>0\).

The axial Hessian block in the gauge is perturbed by \(O(\eta)\) in the row-sum norm. In the diagonal
Schur complement, \(I+E_d\), the same norm of \(E_d\) is \(O(\rho^2\epsilon_A+\rho/j)<1/2\) upon reducing
\(\eta\). Each diagonal incidence row contributes mass \(\le\epsilon_A\). The sum over the board
contributes \(2j\), so the absolute sum of the entries costs \(O(\rho^2j)\). Thus the norm of
\(|\operatorname{tr}E_d^k|\) is bounded by that sum times \(\|E_d\|_\infty^{k-1}\). There are only \(2j\)
axial indices. The action uses \(\sum p^{\rm T}_i|\log p^{\rm T}_i|\le j\log j\) by row entropy, and for
the slacks and local factors it suffices to sum the loads. This gives \(\|\Lambda\|_\rho\le
C(j\log(ej)+\rho^2j)\), proving the stated bound.
\end{proof}

\begin{lema}[Completability on short arrival]\label{lem:positividad-llegada}
For every admissible cut-off and \(n\ge n_0(b)\) we have \(Q(B_j)>0\) on \(H_j\). In particular,
\(H_j\subseteq F_b\).
\end{lema}
\begin{proof}
The scaling \(D\) of the mask has unit margins. For any set of rows its mass must fit in the neighboring
columns. Hall's theorem gives \(\operatorname{per}A>0\). On \(H_j\), \(\rho^2(\epsilon_A+1/j)\le\eta\) by
the choice of \(\alpha\) and \(n_0(b)\). The evaluation at \(Y=A\) of~\eqref{eq:positividad-cubo} gives
\(Z(A)>D_\xi(A)/2>0\). Each rook placement on \(A\) that avoids all diagonal attacks is a completion, and
conversely. Hence \(Q(B_j)=\operatorname{per}A\,Z(A)>0\). The history remained regular up to and
including \(j\), which is the other condition of \(F_b\).
\end{proof}
